\scanpage{8}
Moreover, recalling the definition of
\[
 J_n=\widetilde D_n=\text{set of subobjects of }D_n
\]
we have, for \(\unsure{\mathrm{H}|\forall} n\in\mathbb N\), a canonical, strictly increasing map
\[
 J_n\xrightarrow{\tau_n}\mathbb N \qquad\qquad
 \text{\(\tau_n\), ``type map}
\]
whose image is formed \deleted{of the} of the set
\deleted{\(I_{\leq n}\) of the} \tikz[baseline=.2ex]{\fill (0,0) circle (.025);}% detached point; no punctuation function assigned
\(\mathbb N_{\leq n}=\{n'\in\mathbb N\mid n'\leq n\}\).

Note that each \(J_n\) has a greatest element, which is ``\(D_n\) (as \inserted{the} maximal subobject of \(D_n\)), and \unsure{which I|which \deleted{is} we|which \deleted{he} we|which \deleted{is} we there|which \deleted{he} we there} denote \(d_n\). We therefore have \inserted{(by (ii))}
\[
 \tau_n^{-1}(n)=\{d_n\}
\]
\deleted{If \(n'\leq n\), then for every \(x\in J_n\)}

Let \(n,n'\in\mathbb N\), \deleted{\(\unsure{x\in J|x\in}\)}\deleted{\(x\in J_n,\ n'=\tau_n(x)\), we} \inserted{with \(n'\leq n\)}
consider \deleted{\unsure{via \(x\)|via \((x)\)}} \deleted{we then have \unsure{a|a} map}
\[
 \begin{aligned}
 \mathrm{Hom}(\deleted{\unsure{O|D}}D_{n'},D_n)&\hookrightarrow \widetilde D_n=J_n\\
 u&\mkern2mu\longmapsto u(d_{n'})
 \end{aligned}
\]
which is a bijection from the \(\mathrm{Hom}\) onto the set \(\tau_n^{-1}(n')\), and consider
\[
 \mathrm{Hom}(\deleted{\unsure{D|O}}D_{n'}^{\unsure{4|\ast}},D_n)
 \longrightarrow
 \mathrm{Hom}\bigl(\underset{\substack{\tikz{\draw(0,.08)--(0,0)--(.60,0)--(.60,.08);}\\ J_{n'}}}{\widetilde D_{n'}},\underset{\substack{\tikz{\draw(0,.08)--(0,0)--(.54,0)--(.54,.08);}\\J_n}}{\widetilde D_n}\bigr)
 \quad\raisebox{-3.5ex}{\tikz{\fill (0,.12) ellipse (.025 and .018);\fill (-.07,0) circle (.022);}}% detached points; no punctuation function assigned
\]
we deduce from this a \deleted{bijection} map
\[
 \begin{aligned}
 % One uncertain physical relation: rotated simeq | rotated equals; literal strokes retained.
 \underset{\left[\begin{gathered}
   \unsure{\tikz[x=1em,y=1em,baseline=0pt]{
     \draw[line width=.35pt] (.12,.56) .. controls (.18,.37) and (-.08,.12) .. (.02,.03)
       .. controls (.07,-.02) and (.18,-.02) .. (.23,.03);
     \draw[line width=.35pt] (.43,.59)--(.43,.08);
     \fill (0,.61) circle (.025);}}\\
   \mathrm{Hom}(D_{n'},D_n)
 \end{gathered}\right]}{\tau_n^{-1}(n')}
 &\overset{\mathrm{def}}{=}
 \underset{\substack{\tikz{\draw(0,.07)--(0,0)--(.86,0)--(.86,.07);}\\\text{set of}\\\text{elements}\\\text{of \(J_n\) of}\\\text{type \(n'\)}}}{J_{n',n}\deleted{\unsure{(n')|(n)}}}
 \xrightarrow{\rho_{n,n'}}\mathrm{Hom}(J_{n'},J_n)\\
 &\hspace{7em}\text{or}\quad \xi\longmapsto u_{\xi}^{n}
 \end{aligned}
\]
which is injective, more precisely we have the \unsure{commutativity|commutativity}
\diagram{\begin{tikzcd}[ampersand replacement=\&,column sep=large,execute at end picture={
 \draw ([xshift=.15em,yshift=-1em]\tikzcdmatrixname-1-1.south east)
 ..controls +(-1.1em,.6em) and +(-2.1em,0)..
 node[pos=.95,above,inner sep=.5pt,fill=white] {\(\scriptstyle\text{inclusion}\)}
 ([xshift=.45em,yshift=-1.45em]\tikzcdmatrixname-1-2.south)
 ..controls +(1.9em,-.04em) and +(-1.0em,-.60em)..([xshift=-1em,yshift=-.75em]\tikzcdmatrixname-1-3.south west);}]
 J_{n,n'} \arrow[r,"\rho_{n,n'}"]
 \& \mathrm{Hom}(J_{n'},J_n) \arrow[r,"u\mapsto u(d_{n'})"] \& J_n
\end{tikzcd}}
% archive sheet number: 7
